\section{Introduction}
Mandatory reporting programmes are widespread, and the administrative datasets they generate have become primary sources for empirical research in energy policy, environmental economics, public health, labour economics, and related fields (\cit{liangzhang2026}{Liang \& Zhang, 2026}). Such administrative data is increasingly recognised as a distinctive resource for social science, prized for its scale and population coverage but carrying selection and data-quality challenges that demand explicit methodological treatment (\cit{connellyetal2016}{Connelly et al., 2016}). The U.S. Securities and Exchange Commission requires publicly traded firms to file annual financial disclosures on Form 10-K via the EDGAR system (\cit{securitiesnd}{U.S. Securities and Exchange Commission, n.d.}). The Occupational Safety and Health Administration mandates injury and illness recordkeeping for entities above explicit employee-size thresholds, and requires larger establishments to submit these records electronically (\cit{departmentnd}{U.S. Department of Labor, Occupational Safety and Health Administration, n.d.}). The Environmental Protection Agency's Toxics Release Inventory requires annual reporting of chemical releases and waste management from facilities that meet industry, employment, and chemical-volume thresholds (\cit{environmentalnd}{U.S. Environmental Protection Agency, n.d.}). Numerous municipal and state energy benchmarking ordinances compel building owners to report annual energy consumption, applying to buildings that typically exceed 25,000--50,000 gross square feet (\cit{mengetal2017}{Meng et al., 2017}). In each case, the reporting obligation is triggered by an observable, verifiable characteristic of the reporting unit: firm size, establishment headcount, production or chemical-use volume, or building floor area. This design reflects a regulatory balance between compliance costs and informational benefits.

These programmes share a structural feature that is rarely made explicit: the observable characteristic that triggers the reporting obligation is also a causal determinant of the outcome variable that researchers typically study. A large office building, for example, must report its energy use precisely because it is large, and it emits substantially more greenhouse gases because it is large. Consequently, the compliance threshold does not select units into the observed sample at random. Instead, it systematically enriches the sample with units that exhibit the highest values of the outcome variable. This is not random missingness or dropout; it is systematic unit-level selection on a variable that is simultaneously the selection criterion and a cause of the outcome (\cit{hernanrobins2020}{Hernán \& Robins, 2020, Chapter 8}).

Standard analyses do not correct for this selection mechanism. Regression, gradient-boosted trees, and Bayesian models alike treat the observed sample as representative of the target population unless a correction is added explicitly, and none of the standard modelling pipelines includes one by default. A model fitted to the observed sample estimates the outcome distribution conditional on compliance, P(Y | X, C = 1), when the question of interest is the population relationship P(Y | X). The two quantities differ whenever compliance C is endogenous to the predictors X. This is not a limitation of predictive modelling that can be overcome with more flexible algorithms or more data. It is a causal inference problem: because the sampling mechanism is endogenous to the predictors, associational approaches cannot recover the target population relationship without an explicit model of the selection process.

\cit{pearl2018}{Pearl (2018)} states the underlying limitation directly: no volume of data can answer a causal question without a model of the process that generated it. The failure of conventional workflows is therefore invisible from within those workflows. Cross-validation accuracy does not reveal the bias, because the held-out test set is drawn from the same compliance-enriched distribution as the training set; the model is evaluated against the very population it was fitted to. As a result, high predictive performance on observed data can coexist with systematic distortion of the population-level relationships that researchers and policymakers actually wish to understand.

The ingredients of a solution exist across separate literatures, yet no widely adopted procedure currently integrates them for compliance-gated administrative data. \cit{pearl2019}{Pearl (2019)} identifies seven foundational tools of causal inference and emphasises that causal structure must precede data analysis. However, this framework remains conceptual and does not provide a concrete procedure for recovering that structure from a specific administrative dataset. The causal inference literature offers inverse probability weighting (IPW) as a standard correction for selection bias (\cit{hernanrobins2020}{Hernán \& Robins, 2020}; \cit{rosenbaumrubin1983}{Rosenbaum \& Rubin, 1983}), yet it typically treats the validity of the adjustment set as an untested assumption rather than a property to be tested against the data structure itself. Econometrics has addressed non-randomly selected samples since \cit{heckman1979}{Heckman (1979)}, whose two-step estimator treats selection bias as an omitted-variable specification error; that tradition targets selection driven by unobservables correlated with the outcome, and in its classical form buys identification with distributional or exclusion-restriction assumptions that compliance-gated data, where the selection rule is written down and deterministic in observed covariates, does not require (Section 2.2). Moreover, neither literature specifically addresses compliance-gated data as a distinct class requiring tailored identification strategies.

The machine-learning literature supplies powerful predictive tools but treats the data-generating process as unknown by design, and therefore cannot detect the selection bias introduced by the compliance gate. Prior work on building energy benchmarking illustrates this limitation concretely. \cit{kontokostatull2017}{Kontokosta and Tull (2017)} develop a data-driven framework for predicting building energy use from the New York City benchmarking dataset and achieve strong out-of-sample accuracy using gradient-boosted trees. Similarly, \cit{papadopouloskontokosta2019}{Papadopoulos and Kontokosta (2019)} propose a method to grade buildings along an energy performance curve derived from mandatory disclosure data. Both studies treat the observed benchmarking records as representative of the broader building stock. Neither accounts for the fact that the compliance threshold renders the sample non-representative; as a result, their estimates describe only the compliant subset rather than the full population of buildings. \cit{facure2023}{Facure (2023)} provides practical guidance for applying causal reasoning within machine learning workflows, but treats the causal graph as a modelling input rather than an object of inquiry, offering no procedure for deriving or testing that graph against the data structure of an administrative dataset.

More broadly, the tension between structural and algorithmic modelling approaches (\cit{breiman2001}{Breiman, 2001}) remains unresolved in the presence of selection bias: when both modelling cultures are applied to the same compliance-enriched sample, both estimate the biased quantity P(Y | X, C = 1), and their convergence or divergence carries no information about the target population relationship P(Y | X). Any meaningful cross-paradigm comparison requires a common causal foundation, one grounded in a data-generating process that has been subjected to a falsification test.

The gap is therefore specific and bridgeable. No existing procedure combines falsification testing of the data-generating process against primary institutional documentation, IPW correction grounded in the tested structure, and benchmarking of the resulting structural model against an unconstrained machine-learning model.

This paper supplies that procedure. We propose the Tested DGP (TDGP) pipeline and demonstrate it on Chicago's mandatory energy benchmarking programme, a compliance-gated panel of 28,329 building-year records (2014--2023) drawn from 3,852 buildings, of which 21,714 carry an observed greenhouse-gas value and 5,510 are non-compliant with the reporting obligation; the two counts describe different cuts of the panel rather than a partition of it.

This paper makes four contributions.

First, we formalise the compliance-gated dataset as a distinct analytical class. We define a compliance-gated dataset as a dataset generated under three jointly necessary structural conditions:

\begin{itemize}
  \item Unit-level threshold selection: Reporting is mandatory only if observable selection variables exceed fixed thresholds, with the inclusion decision made separately for each individual unit.
  \item Endogenous selection variable: The selection variables also causally influence the outcome of interest.
  \item Atomic reporting requirement: Units must report their full outcome data or none at all; partial reporting of the outcome is not permitted by the programme design.
\end{itemize}
Together, these three conditions produce a distinctive form of unit-level selection bias that is common in administrative and regulatory datasets but has not been formally characterised as a unified class in the literature. We express this class within the framework of structural causal models (\cit{pearl2009}{Pearl, 2009}).

Second, we propose the Tested DGP (TDGP) pipeline, a structured three-stage procedure for establishing a credible causal graph from compliance-gated administrative data. The DAG establishment problem is the fundamental bottleneck in applied causal inference. Pearl's entire framework is conditional on having a credible causal graph: the do-calculus, the backdoor and front-door criteria, and counterfactual reasoning each presuppose one. Yet in most applied work the graph is asserted rather than tested. For compliance-gated datasets, this problem is tractable: the institutional documentation that defines the compliance obligation also encodes the causal structure. The TDGP procedure exploits this property. The three stages are:

\begin{itemize}
  \item Shadow Matrix Analysis, in which we use the shadow matrix to map the patterns of missingness and locate the compliance-gated block of reported fields in the raw data;
  \item DAG Hypothesis Formulation, where we construct a hypothesis about the Data-Generating Process in the form of a Directed Acyclic Graph (DAG), informed by the missingness patterns, the identified gate, and variable descriptions;
  \item DGP Falsification and Refinement, in which we test and iteratively refine the hypothesised DGP against primary institutional documentation and regulatory rules until the candidate graph has survived a falsification test rather than being merely assumed. In the Chicago application, the EPA accounting formula fixes the mediator structure, and the recorded data bear it out: an OLS reconstruction of GHG from the five mediators achieves R$^2$ $\approx$ 1.000. Because the standard \emph{defines} emissions as that weighted sum of metered fuels, this reconstruction confirms that the data conform to a documented accounting identity rather than an empirically discovered relationship, which is precisely what makes the structural claim $\mathrm{pa}(Y) = M$ checkable here.
\end{itemize}
What the pipeline tests and what it must still assume are kept distinct throughout: the documents settle the selection rule, the mediator/outcome classification, and the claim $\mathrm{pa}(Y) = M$, whereas the ignorability of any residual, unobserved selection is reached by no document and is carried as a stated assumption (Sections 3.4 and 8). Once these three stages are complete, the tested DAG licenses the full Pearlian causal cycle of association, intervention, and counterfactual as a direct logical consequence, not an additional methodological claim. We demonstrate this by applying IPW selection correction, hierarchical Bayesian estimation, and counterfactual inference to the Chicago dataset. To the best of our knowledge, this is the first applied study to specify a complete causal graph and subject it to a falsification test against primary institutional documentation in an administrative dataset, and to demonstrate the full Pearlian causal framework on that basis (\cit{pearlmackenzie2018}{Pearl \& Mackenzie, 2018}).

Third, we demonstrate that applying selection bias correction has a deeper benefit than simply removing bias: it creates a common causal foundation that allows models from fundamentally different traditions to answer the same population-level questions. By correcting for the compliance-induced selection mechanism, both structural Bayesian models and machine-learning ensembles can be interpreted as estimating the same causal quantities. This comparability is theoretically justified by the mathematical equivalence of the structural causal model and potential outcomes frameworks (\cit{shpitserpearl2008}{Shpitser \& Pearl, 2008}), which permits the use of Pearl's graphical language for reasoning alongside Rubin's weighting methods for estimation. Furthermore, the work of \cit{bareinboimpearl2012}{Bareinboim and Pearl (2012)} provides the formal foundation for recovering population-level causal effects from non-representative samples. Bias correction does more than improve validity; it enables genuine cross-paradigm comparison within a unified causal framework. Whether the paradigms then agree is an empirical question, taken up as the fourth contribution.

Fourth, we demonstrate empirically that a structural model whose form has been subjected to a falsification test can achieve comparable predictive accuracy to a machine-learning model while delivering substantial additional value that purely algorithmic approaches cannot provide. A hierarchical Bayesian model with a falsification-tested structural form achieves MdAPE = 21.4\% versus 21.0\% for an unconstrained gradient-boosted ensemble, a gap of less than one percentage point. When the data-generating process can be reliably recovered from institutional documentation, the structural model is neither less accurate nor less useful than its algorithmic counterpart; it additionally yields interpretable causal parameters, well-calibrated uncertainty estimates, and full counterfactual capability. This finding is a corollary of the DGP surviving its falsification test: when the causal graph has been tested rather than merely assumed, structural and algorithmic models converge because they are estimating the same quantity. It directly addresses the long-standing tension articulated by \cit{breiman2001}{Breiman (2001)} between structural and algorithmic modelling cultures, and it identifies the precise condition under which that tension dissolves. We demonstrate this on a single programme, so the result is best read as a worked existence proof, that the condition can hold and what follows when it does, rather than a general resolution of the debate; whether the agreement recurs across the class is the empirical question Section 8 leaves open. The TDGP pipeline is applicable to any dataset satisfying the three structural conditions that define compliance-gated data, a class spanning energy, finance, labour, and healthcare reporting systems.

The remainder of the paper is organised as follows. Section 2 formalises the compliance-gated dataset as a distinct analytical class, defines the three jointly necessary structural conditions, and establishes why this class makes the DAG falsification problem tractable. Section 3 presents the Tested DGP (TDGP) pipeline in dataset-agnostic form, describing the Shadow Matrix analysis, DAG hypothesis formulation, and DGP falsification procedure that constitute its three stages, then the causal analysis the tested DAG licenses. Section 4 executes the pipeline on Chicago's mandatory energy benchmarking programme, presenting the shadow matrix results, the constructed DAG, and the EPA formula test that leaves the hypothesised causal structure unrefuted. Section 5 demonstrates the causal estimation that follows from the tested DAG: a 2$\times$3 factorial design across six models isolates the contribution of the causal pipeline from the contribution of functional form, with the model comparison table and the central $\Delta$ MdAPE < 1 pp finding reported here. Section 6 reports the structural inference from the focal model: the posterior distributions over the four parameters of interest and a sensitivity analysis showing robustness to prior specification and IPW trim cap. Section 7 demonstrates counterfactual reasoning as a direct consequence of the tested DAG, evaluating two policy-relevant counterfactuals. Section 8 concludes by restating the core finding, delineating the scope and transferability of the pipeline, acknowledging its limitations, and sketching directions for future work.

